Let $\mathcal{U} = \{u_i\}_{i=1}^{N}$ and $\mathcal{T} = \{t_j\}_{j=1}^{M}$ be
the set of $N$ users and $M$ items of a system, and let
$\mathcal{I} = \{(u_i, t_j)\}$ be the observed implicit interactions, where
$(u_i, t_j) \in \mathcal{I}$ states that user $u_i$ has interacted with item
$t_j$. Indices are omitted whenever a statement holds for an arbitrary user or
item. Matrix factorization assigns a learned embedding $\mathbf{u}$ to every
user and a learned embedding $\mathbf{t}$ to every item and scores a pair by
the dot product of the two. ProtoMF~\cite{melchiorre2022protomf} keeps these
embeddings and adds prototypes, which are further learned vectors placed in
the same space as the entities they describe. The model is built in three
steps, prototypes on the user side, prototypes on the item side, and both at
once, and we follow that order.

U-ProtoMF rests on the assumption that the users of a system share a limited
number of item consumption patterns, each of which can be embodied by a
prototypical user. One user prototype may stand for the preference for Folk
Metal tracks, another for the enjoyment of drama and romance
movies~\cite{melchiorre2022protomf}. The model introduces a set $\mathcal{P}^u$
of $L^u$ user prototypes with $L^u \ll N$, each of them a learned vector
$\mathbf{p}^u_l$ of the embedding dimension $d$. Every user then receives a
second representation $\mathbf{u}^*$, the vector holding the similarity of the
user embedding to each of the prototypes,
\begin{equation}
\mathbf{u}^* =
\begin{pmatrix}
\mathrm{sim}(\mathbf{u}, \mathbf{p}^u_1) \\
\vdots \\
\mathrm{sim}(\mathbf{u}, \mathbf{p}^u_{L^u})
\end{pmatrix}
\in \mathbb{R}^{L^u},
\qquad
\mathrm{sim}(\mathbf{a}, \mathbf{b}) = 1 +
\frac{\mathbf{a}^\top \mathbf{b}}
     {\lVert \mathbf{a} \rVert \cdot \lVert \mathbf{b} \rVert},
\label{eq:proto-sim}
\end{equation}
where $\mathrm{sim}$ is the shifted cosine similarity and
$\lVert \mathbf{x} \rVert$ is the $L^2$ norm of $\mathbf{x}$. The shift keeps
every entry of $\mathbf{u}^*$ positive and inside the range from 0 to 2.

The affinity score of a user and an item is a linear combination of this new
user representation and the item embedding,
\begin{equation}
\text{U-score}(u,t) = \sum_{l=1}^{L^u} s^{\mathrm{user}}_l,
\qquad
\mathbf{s}^{\mathrm{user}} = \mathbf{u}^* \odot \mathbf{t},
\label{eq:u-score}
\end{equation}
where $\odot$ is the element-wise product and the item embedding
$\mathbf{t} \in \mathbb{R}^{L^u}$ now carries one entry per user prototype.
Equation~\ref{eq:u-score} is the dot product
$\text{U-score}(u,t) = \mathbf{u}^{*\top} \mathbf{t}$, written in a form that
exposes the intermediate vector $\mathbf{s}^{\mathrm{user}}$, which breaks the
score down into one contribution per user prototype.

Three properties of this construction carry the model's interpretability
claim. Every coordinate of $\mathbf{u}^*$ belongs to one prototype and
therefore to one consumption characteristic, so a user is described one
characteristic at a time. The item embedding acts as the vector of weights an
item places on those characteristics, and a heavy metal track is expected to
carry a high weight on the coordinate of the prototype that represents metal
fans. The score is a linear function of the similarities, so every prototype
holds a quantifiable share of a given recommendation. Two further points
follow from the shape of $\mathbf{u}^*$. The representation is graded rather
than exclusive, since a user is related to all prototypes at once instead of
being assigned to one of them, which allows several tastes to coexist in one
profile. And that profile is not a description placed beside the model, it is
the representation the score is computed from. Anchoring a latent space with a
small set of shared vectors is not unique to ProtoMF, and anchor-based
collaborative filtering~\cite{barkan2021acf} is its closest relative, using
one shared set of anchors for both users and items.

I-ProtoMF applies the same construction to the item side, where the assumption
is that items are co-consumed in a limited number of patterns, so that an item
prototype represents, for example, the items of one musical genre or product
category. The model learns a set $\mathcal{P}^t$ of $L^t$ item prototypes
$\mathbf{p}^t_l$ with $L^t \ll M$ and represents every item by its
similarities to them,
\begin{equation}
\mathbf{t}^* =
\begin{pmatrix}
\mathrm{sim}(\mathbf{t}, \mathbf{p}^t_1) \\
\vdots \\
\mathrm{sim}(\mathbf{t}, \mathbf{p}^t_{L^t})
\end{pmatrix}
\in \mathbb{R}^{L^t}.
\label{eq:proto-tstar}
\end{equation}
The score mirrors Equation~\ref{eq:u-score},
\begin{equation}
\text{I-score}(u,t) = \sum_{l=1}^{L^t} s^{\mathrm{item}}_l,
\qquad
\mathbf{s}^{\mathrm{item}} = \mathbf{t}^* \odot \mathbf{u},
\label{eq:i-score}
\end{equation}
where the user embedding $\mathbf{u} \in \mathbb{R}^{L^t}$ now carries one
entry per item prototype and the score is again a dot product,
$\text{I-score}(u,t) = \mathbf{u}^\top \mathbf{t}^*$. The element-wise form
serves the same purpose as before and splits the score over the item
prototypes.

Both models are trained on the observed interactions with the
cross-entropy loss over the softmax of the scores,
\begin{equation}
\mathcal{L}_{rec}(\mathcal{I}, \Theta) =
- \sum_{(u,t) \in \mathcal{I}} \ln p(t \mid u) +
\lambda_{L2} \lVert \Theta \rVert,
\qquad
p(t \mid u) = \frac{e^{\,\text{U-score}(u,t)}}
                   {\sum_{j=1}^{M} e^{\,\text{U-score}(u,t_j)}},
\label{eq:proto-rec-loss}
\end{equation}
where $\Theta$ holds the model parameters and $\lambda_{L2}$ weights the $L^2$
regularization term. The loss of I-ProtoMF is the same with
$\text{I-score}$ in place of $\text{U-score}$. ProtoMF adds two further terms
to this objective, which act on the geometry of the prototype space and are
treated in Section~\ref{sec:bg-protomf-reg}.

The prototypes themselves are learned parameters and carry no name of their
own, so both models are read after training, and the two sides are read
differently. Item prototypes share their space with the item embeddings, so an
item prototype is interpreted through its nearest items, and whatever their
metadata has in common becomes its name. A user prototype cannot be read this
way, because the objects a person recognizes are items rather than users. It
is interpreted instead by building a synthetic representation $\mathbf{u}^*$
that assigns the maximum similarity to the prototype in question and zero to
all others, and by reading the recommendations the model returns for that
imaginary user~\cite{melchiorre2022protomf}. In both cases the name of a
prototype is a summary of the neighborhood it sits in, obtained from the
trained model by inspection.
